\section{The exact whole lower index of an original crossing run}

The original crossing lower history retains the same marked datum and a true
interior empty step. Its actual insertion index is its WHOLE original lower
boundary length plus the true remaining head length. Actual SNC geometry
computes that index from the original level; it is not a comparison premise.
This refines the indexed Step One bookkeeping of
\cite{wlodarczyk05-hironaka,bierstone-milman08-functoriality}.

\begin{lemma}[Whole original lower boundary length at the actual level]
\label{hir:original-whole-lower-boundary-exact-length}
\source{wlodarczyk05-hironaka,bierstone-milman08-functoriality}
\uses{hir:original-member-filtered-prefix-exact-count,hir:step-one-bridge}
\lean{Hironaka.HironakaOriginalCrossingExactIndex.original_lower_boundary_length_actual_level}
\leanok
Let $M$ be originally well formed, let $K$ be an actual bar-test history, and
let $p$ be an actual stage point with level $\ell$. The WHOLE actual lower
marked record formed at its true level has boundary length
\[
 |\operatorname{lowerMarkedM}(M,\ell,K,p).E|
       =|M.E|+|K|-\ell.
\]
At level zero this is the genuine chosen-contact lower record; at every
positive level it is the genuine member-restricted lower record. No carrier,
choice, index or boundary-count certificate is assumed.
\end{lemma}
\begin{proof}
At level zero the actual coefficient/contact record retains the entire bar
transform boundary, and open pullback preserves its length. The true bar
transform length is $|M.E|+|K|$, and subtraction by zero leaves it unchanged.
At positive level the actual lower boundary is the open pullback of the
entire bar boundary filtered by removal of the original members through $p$.
Actual bar-testness makes this bar transform genuinely well formed and SNC.
Its literal original member list is its prefix of length $|M.E|$. Apply the
proved exact original filtered-prefix count with larger prefix equal to its
WHOLE boundary. The boundary count of the original prefix is the true level
$\ell$. The map into the actual open lower carrier preserves the retained
length, yielding the displayed identity. This counts every actual boundary
index and permits inactive or empty repeated entries.
\end{proof}

\begin{theorem}[Original crossing canonicity internally computes the exact lower index]
\label{hir:original-crossing-lower-index-exact}
\source{wlodarczyk05-hironaka,bierstone-milman08-functoriality}
\uses{hir:original-whole-lower-boundary-exact-length,hir:original-canonical-crossing-lower-data,hir:original-interior-crossing-selected-splits}
\lean{Hironaka.HironakaOriginalCrossingExactIndex.original_canonical_crossing_lower_index_exact}
\leanok
Assume only an original dimension-$(k+1)$ maximal-order datum and an actual
canonical Step One history $HS$ for the specific original lower family
$R_k$. Let $y$ be an actual nonspecial point whose TRUE own-level run
$r$ reaches the insertion after $H$, so $r\geq|S|$. Put
$K=\operatorname{head}(H,r-|S|)$, $B=\operatorname{tail}(H,r-|S|)$, $p=(BS)(y)$,
$\ell=\operatorname{level}(M,HS,y)$, and
$J=\operatorname{lowerMarkedM}(M,\ell,K,p)$.
Then the exact actual lower insertion position is
\[
 |J.E|+|B|=|M.E|+|H|-\ell.
\]
Every actual lower order, including positive recursive orders, is allowed.
\end{theorem}
\begin{proof}
The genuine original canonical bridge supplies actual whole bar-testness and
hence bar-testness of the literal selected prefix $K$. Its actual nonspecial
current point lies in its own level regime. Since the true suffix length
$|B|+|S|$ is at most $r$, actual reaching-run membership gives the SAME level
$\ell$ at the assignment point $p$. Apply the preceding whole lower-length
identity there. This also gives $\ell\leq|M.E|+|K|$, since the true level counts
indices in an original prefix. Actual head/tail lengths sum to $|H|$.
Consequently adding $|B|$ to the computed whole lower boundary length yields
exactly $|M.E|+|H|-\ell$. Nothing about lower values or inserted canonical
histories is assumed.
\end{proof}

\begin{theorem}[The actual crossing leaf uses the same whole Step One history map]
\label{hir:original-crossing-leaf-common-history-map}
\source{wlodarczyk05-hironaka,bierstone-milman08-functoriality}
\uses{hir:original-crossing-lower-index-exact,hir:original-canonical-crossing-step-one-monomial-leaf,hir:original-short-history-step-map}
\lean{Hironaka.HironakaOriginalCrossingExactIndex.original_canonical_crossing_leaf_history_map}
\leanok
For the preceding original canonical inputs AT AN ACTUAL ORIGINAL COSUPPORT
POINT, assume the suffix $S$ is nonempty.
The genuine original lower-data producer internally constructs the whole
actual lower canonical tail, domain and cosupport point. On its actual lower
nonmonomial-order-zero branch, the ENTIRE original Step One value on
$H\epsilon S$ is $F_{k,|M.E|+|H|}$ of the original value on $HS$.
The original point may have zero or positive Step One level.
\end{theorem}
\begin{proof}
The genuine original canonical crossing-leaf producer supplies the actual
lower witness, domain, full canonical tail and cosupport internally. On its
stated actual order-zero branch it proves the complete Step One pair formula
with actual family index $|J.E|+|B|$. Substitute the exact original index
computed above. The first coordinate of the whole original Step One value
is its actual level, so this pair formula is precisely the definition of the
single level-indexed strict history map $F_{k,j}$. The lower order-zero
condition is an explicit real branch on internally supplied original data;
positive crossing recursion remains open. Together with the complete short
Step One formula, this supplies the common map for actual short points and
actual crossing leaf points, without claiming full higher crossing transport
or noncompact scheduling.
\end{proof}
